\par\medskip
\hypertarget{native-070-article053}{}\pdfbookmark[2]{53.}{article-053}
\begin{center}53.\end{center}
The position of any point in space is most conveniently determined by its distances from three planes that intersect at right angles. Taking the plane of the ecliptic as one of the planes and denoting by $z$ the distance of the celestial body from it, counted as positive on the north side and negative on the south side, we clearly have
\[
z=r'\operatorname{tang}\beta=r\sin\beta=r\sin i\sin u.
\]
The other two planes, which we also assume to pass through the Sun, project great circles onto the celestial sphere. These cut the ecliptic at right angles, so their poles lie on the ecliptic itself and are $90^0$ apart. For either plane, we call the pole on whose side distances are counted as positive the \textit{positive pole}. Let $N$ and $N+90^0$ be the longitudes of the positive poles, and denote the distances from the corresponding planes by $x$ and $y$, respectively. \SourceLeaf{071}It is then easy to see that
\[
\begin{aligned}
x&=r'\cos(\lambda-N)\\
&=r\cos\beta\cos(\lambda-\AscendingNode)\cos(N-\AscendingNode)\\
&\quad+r\cos\beta\sin(\lambda-\AscendingNode)\sin(N-\AscendingNode)\\[.4em]
y&=r'\sin(\lambda-N)\\
&=r\cos\beta\sin(\lambda-\AscendingNode)\cos(N-\AscendingNode)\\
&\quad-r\cos\beta\cos(\lambda-\AscendingNode)\sin(N-\AscendingNode),
\end{aligned}
\]
which become
\[
\begin{aligned}
x&=r\cos(N-\AscendingNode)\cos u+r\cos i\sin(N-\AscendingNode)\sin u\\
y&=r\cos i\cos(N-\AscendingNode)\sin u-r\sin(N-\AscendingNode)\cos u.
\end{aligned}
\]
Thus, if the positive pole of the plane corresponding to $x$ is placed at the ascending node itself, so that $N=\AscendingNode$, we obtain the simplest expressions for the coordinates $x$, $y$, and $z$
\[
\begin{aligned}
x&=r\cos u\\
y&=r\cos i\sin u\\
z&=r\sin i\sin u.
\end{aligned}
\]
If this assumption does not hold, the formulas given above nevertheless take an almost equally convenient form when we introduce four auxiliary quantities $a$, $b$, $A$, and $B$, determined so that\HistoricalNoteEight{Manuscript note, for which see p. 69, line 7:]
\[
\begin{aligned}
-\cos\AscendingNode'&=a\sin A\\
\cos i'\sin\AscendingNode'&=a\cos A\\
-\sin\AscendingNode'&=b\sin B\\
-\cos i'\cos\AscendingNode'&=b\cos B.
\end{aligned}
\]
}
\[
\begin{aligned}
\cos(N-\AscendingNode)&=a\sin A\\
\cos i\sin(N-\AscendingNode)&=a\cos A\\
-\sin(N-\AscendingNode)&=b\sin B\\
\cos i\cos(N-\AscendingNode)&=b\cos B
\end{aligned}
\]
(see Article 14, II). It is then clear that
\[
\begin{aligned}
x&=ra\sin(u+A)\\
y&=rb\sin(u+B)\\
z&=r\sin i\sin u.
\end{aligned}
\]
